% !TEX encoding = UTF-8 Unicode
\documentclass[12pt,a4paper]{report}
 
\usepackage[brazil]{babel}
\usepackage[utf8]{inputenc}
\usepackage[T1]{fontenc}
\usepackage{graphicx, subfigure}
\usepackage{epstopdf}
\usepackage{indentfirst}
\usepackage{fancyhdr}

\graphicspath{ {../images/} }

\epstopdfsetup{outdir=./}
\epstopdfsetup{suffix=}

\pagestyle{fancy}

\fancyhead[RO,RE]{\includegraphics{unb.eps}}
\fancyhead[LO,LE]{Departamento de Ciência da Computação\\
116441 --- Engenharia de Software}

\fancypagestyle{plain}{
\fancyhead[RO,RE]{\includegraphics{unb.eps}}
\fancyhead[LO,LE]{Departamento de Ciência da Computação\\
116441 --- Engenharia de Software}}

\title{Plano de Projeto do Sistema de Gestão Eletrônica de Documentos da Defensoria Pública do Distrito Federal - SGED/DPDF}
\author{Pedro Salum\\
	09/0139232\\
	pedro@loopec.com.br
	\and
	Daniel Sandoval\\
	09/0109899\\
	daniel@loopec.com.br}
	
\begin{document}
\maketitle
\tableofcontents

\chapter{Introdução}

Este documento tem por objetivo explicitar o planejamento para a realização do projeto de Gestão Eletrônica de Documentos da Defensoria Pública do Distrito Federal --- SGED/DPDF --- no tocante ao desenvolvimento, arquitetura, testes, implantação, manutenção e riscos. A seguir há uma breve descrição de cada capítulo.

\paragraph{Processo de Desenvolvimento} Discorre sobre decisões tomadas em relação a metodologias de desenvolvimento de \textit{software} e gerência de projeto. Discute a escolha de metodologias ágeis sobre métodos planejados.

\paragraph{Arquitetura e padrões de projeto} Discorre sobre escolhas de arquitetura e padrões de projeto que serão utilizados na codificação do sistema.

\paragraph{Plano de Testes} Discorre sobre as metodologias e ferramentas de testes escolhidas para o projeto, considerando deste testes de código até de aceitação, como e quando serão realizados.

\paragraph{Plano de Implantação e Manutenção} Discorre sobre as metodologias e ferramentas escolhidas para o processo de implantação e manutenção contínua do sistema entregue.

\paragraph{Análise de Riscos} Discorre sobre riscos que devem ser considerados ao longo do projeto e possibilidades de mitigação.

\chapter{Processo de Desenvolvimento}

A metodologia de desenvolvimento escolhida foi baseada em métodos ágeis. Levando-se em consideração que o sistema a ser desenvolvido está inserido em um ambiente corporativo, temos que estar preparados para mudanças. Mudanças de prioridades; regras de negócio; gerência; cronograma; entre outras, devem ser suportadas com o mínimo de prejuízo possível. Além disso, há experiências bem sucedidas com desenvolvimento ágil em ambientes corporativos governamentais \cite{metaginstpub}.

Sabendo que métodos ágeis não são próprios para implantação uniforme em qualquer situação \cite{cohn2009succeeding}, as ferramentas e metodologias escolhidas para este projeto são uma sugestão e podem sofrer modificações ou adaptações ao longo do desenvolvimento.

\paragraph{Scrum} 
%-> Definição e explicação de qual processo de desenvolvimento de software escolhido para o projeto

\chapter{Arquitetura e padrões de projeto}

%-> Definição e explicação da arquitetura de software escolhida para o projeto, focalizando a escolha de padrões de arquitetura, assim como diagrama de arquitetura
%-> Definição e explicação dos design patterns escolhidos para serem implementados no projeto

\chapter{Plano de Testes}
%-> Discussão dos testes de verificação e validação realizados no projeto, assim como descrição de casos de testes, inspeção de código, testes unitários

\chapter{Plano de Implantação e Manutenção}
\paragraph{Implantação} A implantação do sistema aqui descrito será por meio de um deploy em um servidor central da Defensoria Pública do Distrito Federal e uma cópia do mesmo em cada um dos 26 Núcleos de Assistência Jurídica espalhados pelo Distrito Federal. A cópia dentro do servidor central será configurada como "MASTER", ou seja, recebe informações das outras 26 cópias espalhadas, configuradas como "SLAVE", e as armazena para o próximo backup.

\paragraph{Manutenção} Manutensões serão efetuadas em ambiente de teste, as quais após estarem estáveis e prontas para deploy, seráo mandadas para o servidor central, substituindo a cópia "MASTER", onde os "SLAVES" irão solicitar a versão mais nova do sistema o mais cedo possível, ou seja, uma cópia do "MASTER" será baixada para o "SLAVE" a qual substituirá a versão anterior. Importante ressaltar que só serão transmitidos as partes atualizadas do sistema, reduzindo a carga da rede.

\chapter{Análise de Riscos}

A pesar de ser um software consideravelmente pequeno, quando falamos se sistemas coorporativos, ainda existem riscos que devem ser levados em consideração, mesmo que estes sejam completamente improváveis de acontecerem no curto-médio prazo. A seguir estão listados, em nível de probabilidade, e detalhados os riscos aqui levados em consideração.

\paragraph{Alta} Eventos de álta probabilidade são aqueles referentes à equipe de desenvolvimento e suas relações com as partes, tanto UnB quanto DP-DF. Dentre estes riscos estão:

\begin{itemize}
\item [-] Saída de membro da gerência do projeto
\item [-] Saída de membro altamente qualificado, tanto na área de teste quanto de desenvolvimento
\item [-] Atraso no recebimento do capital para pagamento da equipe
\item [-] Falta de cooperação entre a Defensoria Públia do Distrito Federal e a equipe interna.
\end{itemize}

Para evitarmos problemas internos da equipe, iremos bem remunerá-los, manter reservas financeiras proveninentes do não pagamento dos coordenadores no início do projeto, remunerando-os no final, gerando um caixa de segurança. Além da perte financeira, iremos prover toda a estrutura necessária para que o trabalho possa ser efetuado além de ambientes para descanso e descontração. Para evitar a falta de cooperação da equipe interna da DP-DF, iremos negociar com a diretoria do órgão para que um membro seja alocado somente para se comunicar com a nossa equipe, retirando os outros afazeres do mesmo.

\paragraph{Média} Eventos de média probabilidade são aqueles internos à instituiçãoe que são possíveis de ocorrer e que podem impedir ou dificultar o desenvolvimento do sistema como um todo. Estes são:
\begin{itemize}
\item [-] Cancelamento do Contrato entre DP-DF e a Universidade de Brasília
\item [-] Troca do quadro de diretores
\item [-] Probelmas de infraestrutura na área de TI, sejam com servidores ou com a parte de telecomunicação
\end{itemize}

Visando a redução destes riscos de média probabilidade, devemos efetuar a prestação de serviço de desenvolvimento de software o mais rápido possível, reduzindo os riscos de existir uma possível troca de gestão e cancelamento de contrato entre as partes.
Sobre problemas de infraestrutura de TI, o sistema deverá funcionar tanto de forma local (Por meio de uma LAN) quanto em rede (WAN). Sendo assim, devemos partir para uma abordagem distribuida, com possível topologia de Malha ou Estrela extendida.

\paragraph{Baixa} Eventos de baixa probabilidade são aqueles que modificam estruturas que são externas à Defensoria Pública do Distrito Federal, e não cabe à instituição decidi-las. Dentre elas estão:
\begin{itemize}
\item [-] Modificação na estrutura judiciária do Distrito Federal
\item [-] Modificação de como ocorre a comunicação entre as partes (DP-DF - Tribunal)
\end{itemize}

Com o objetivo de reduzir possíveis danos gerados por estes riscos, o sistema deve ser flexível o suficiente para se adaptar à tais mudanças, seja por meio de poucas modificações ou, idealmente, nenhuma modificação.

\paragraph{Muito Baixa} Eventos de probabilidade baixa podem ser considerados catastróficos, ou seja, não só inviabilizam o desenvolvimento do software como impossibilitam o seu uso e existência da própria instituição. Entre eles estão:
\begin{itemize}
\item [-] Terremoto com epicentro próximo à Brasília devasta toda a estrutura jurídica de Brasília e todo o recurso é devidamente realocado
\item [-] Meteoro atinge o prédio da DP-DF
\item [-] Godzilla surge do fundo do Lago Paranoá e derruba o prédio da DP-DF
\item [-] ETs invadem a terra e sugam o prédio da DP-DF
\end{itemize}


Devido a baixa probabilidade de ocorrencia dos riscos de importância Muito Baixa, e são impossíveis de serem mitigados, estes serão descartados para fins deste projeto. 

%-> Discussão dos riscos do projeto e estratégia de mitigação
\bibliographystyle{plain}
\bibliography{es}

\end{document}
